% !TeX root = ../main.tex

% ==================================================
% 2.4. 인스턴스 속성과 네임스페이스 — self.name은 어디에 저장되는가?
% ==================================================

\subsection{인스턴스 속성과 네임스페이스 - \zsh{self.name}은 어디에 저장되는가?}

2.2에서는 \zsh{class}문을 실행하면 클래스 객체가 생성된다는 사실을 학습했고,

2.3에서는 클래스 객체와 인스턴스 객체가 서로 다른 객체라는 사실을 확인했다.

이번 2.4에서는 한 단계 더 나아가 인스턴스가 자신의 데이터를 어떻게 보관하고,

Python이 그 데이터에 어떻게 접근하는지 살펴보자.

특히 다음 코드에 집중할 것이다.

\begin{dhkcode}
class Person:
    def __init__(self, name):
        self.name = name

p = Person("Kim")
\end{dhkcode}

우리는 일반적으로 \zsh{self.name = name}을 인스턴스에 속성을 저장하는 코드라고 설명한다.

그런데 여기서 몇 가지 질문이 생긴다.

\begin{dhksmry}
1. \zsh{self.name}에서 \zsh{self}와 \zsh{name}은 각각 무엇인가?

2. \zsh{name}이라는 속성은 실제로 어디에 저장되는가?

3. \zsh{p.name}과 \zsh{name}은 왜 다른가?

4. \zsh{p1}과 \zsh{p2}가 서로 다른 값을 저장할 수 있는 이유는 무엇인가?

5. \zsh{p.name}을 읽을 때 Python은 어디에서 값을 찾는가?

6. 클래스 속성과 인스턴스 속성이 같은 이름을 가지면 어떻게 되는가?

7. 이 구조가 PyTorch의 \zsh{self.linear}와 어떻게 연결되는가?
\end{dhksmry}

이번 학습의 핵심은 이름, 객체, 네임스페이스, 속성 접근을 하나의 구조로 연결하는 것이다.

\subsubsection{개념 정의 - 속성과 네임스페이스}

먼저 두 용어를 구분한다.

\zsh{속성(attribute)}은 객체에 연결된 이름으로, 점(.)을 사용하는 속성 접근 문법을 통해 읽거나 설정할 수 있는 값이다.

\zsh{네임스페이스(Namespace)}는 이름과 객체 사이의 바인딩을 관리하는 공간이다.

예를 들어 다음 코드를 보자.

\begin{dhkcode}
name = "Kim"
\end{dhkcode}

여기서 \zsh{name}은 현재 네임스페이스에 존재하는 이름이다.

반면 다음 코드는 다르다.

\begin{dhkcode}
p.name = "Kim"
\end{dhkcode}

여기서 \zsh{name}은 객체 \zsh{p}의 속성 이름이다.

두 표현식은 모두 \zsh{"Kim"}이라는 문자열 객체를 참조할 수 있지만, 이름을 관리하는 위치와 접근 방법이 다르다.

\begin{dhktable}
|표현식|	의미|	이름을 찾는 방식|
|:---|:---|:---|
|name|	현재 실행 환경의 일반적인 이름|	스코프 및 이름 탐색 규칙|
|p.name|	p 객체의 name 속성|	객체의 속성 접근 규칙|
|Person.name|	Person 클래스 객체의 name 속성|	클래스 객체의 속성 접근 규칙|
\end{dhktable}

따라서 \zsh{name}과 \zsh{p.name}은 같은 이름을 공유하더라도 별개의 바인딩일 수 있다.

직접 확인해 보자.

\begin{dhkcode}
name = "Lee"

class Person:
    pass

p = Person()
p.name = "Kim"

print(name)
print(p.name)
\end{dhkcode}

실행 결과:

\begin{dhkzsh}
Lee
Kim
\end{dhkzsh}

왜 다른 값이 출력될까?

\zsh{name}은 현재 모듈의 전역 네임스페이스에 있는 이름이고,

\zsh{p.name}은 인스턴스 객체에 연결된 속성이기 때문이다.

두 이름은 서로 다른 공간에서 관리된다.

\subsubsection{인스턴스 속성은 어디에 저장되는가? - \zsh{__dict__}}

Python의 일반적인 사용자 정의 클래스 인스턴스는

속성들을 \zsh{__dict__}라는 딕셔너리에 저장한다.

다음 코드를 보자.

\begin{dhkcode}
class Person:
    pass

p = Person()

p.name = "Kim"
p.age = 30

print(p.__dict__)
\end{dhkcode}

실행 결과:

\begin{dhkzsh}
{'name': 'Kim', 'age': 30}
\end{dhkzsh}

여기서 중요한 사실을 알 수 있다.

우리는 다음처럼 속성을 설정했다.

\begin{dhkcode}
p.name = "Kim"
p.age = 30
\end{dhkcode}

그런데 인스턴스의 \zsh{__dict__}를 확인하니 이름과 값이 딕셔너리 형태로 관리되고 있다.

개념적으로 다음과 같은 구조다.

\begin{dhkgraph}
node p (1,1) :: p
node d (1,2) [accent] :: __dict__ //"name" → "Kim" // "age" → 30

group per(d) :: Person 인스턴스

edge p -> per

caption :: 일반적인 사용자 정의 인스턴스의 속성 저장 구조를 단순화한 그림
\end{dhkgraph}

여기서 \zsh{p}는 인스턴스 객체를 참조하는 이름이고, 인스턴스는 자신의 속성 정보를 \zsh{__dict__}에 보관한다.

따라서 일반적인 경우 다음 두 표현식으로 같은 값을 확인할 수 있다.

\begin{dhkcode}
print(p.name)
print(p.__dict__["name"])
\end{dhkcode}

실행 결과:

\begin{dhkzsh}
Kim
Kim
\end{dhkzsh}

하지만 두 표현식이 내부적으로 항상 동일하게 처리되는 것은 아니다.

\zsh{p.name}은 Python의 속성 접근 규칙을 따르고, \zsh{p.__dict__["name"]}은 딕셔너리에 직접 접근한다.

이 차이는 이후 속성 탐색과 디스크립터를 이해할 때 중요해진다.

또한 모든 객체가 \zsh{__dict__}를 가지는 것은 아니다.

예를 들어 \zsh{__slots__}를 사용하는 클래스에서는 인스턴스 속성을 다른 방식으로 관리할 수 있다.

현재는 일반적인 사용자 정의 클래스에서 인스턴스 속성이 \zsh{__dict__}에 저장된다는 원리를 이해하면 충분하다.

\subsubsection{\zsh{self.name = name}의 정확한 의미}

이제 처음의 코드로 돌아가 보자.

\begin{dhkcode}
class Person:
    def __init__(self, name):
        self.name = name
\end{dhkcode}

이 코드에서 \zsh{self.name = name}을 이해하려면 왼쪽과 오른쪽을 구분해야 한다.

\begin{dhkcode}
self.name = name
#  왼쪽      오른쪽
\end{dhkcode}

\begin{dhknote}
\begin{dhksmry}[centered]
오른쪽 name
\end{dhksmry}
\end{dhknote}

\zsh{__init__()} 함수의 매개변수 이름이다.

함수가 호출될 때 전달된 객체를 참조하도록 바인딩된다.

\begin{dhknote}
\begin{dhksmry}[centered]
왼쪽 self
\end{dhksmry}
\end{dhknote}

현재 초기화 중인 인스턴스 객체를 참조하는 매개변수다.

\begin{dhknote}
\begin{dhksmry}[centered]
왼쪽 self.name
\end{dhksmry}
\end{dhknote}

\zsh{self}가 참조하는 인스턴스에 \zsh{name}이라는 속성을 설정하라는 의미다.

예를 들어 다음 코드를 실행한다고 생각해 보자.

\begin{dhkcode}
p = Person("Kim")
\end{dhkcode}

일반적인 인스턴스 생성 과정에서 \zsh{__init__()}이 호출될 때는 개념적으로 다음과 같은 인수 전달이 이루어진다.

\begin{dhkcode}
Person.__init__(새로_생성된_인스턴스, "Kim")
\end{dhkcode}

따라서 \zsh{__init__()} 내부에서는 다음과 같은 관계가 성립한다.

\begin{dhksmry}
    self \ra 새로 생성된 Person 인스턴스

    name \ra "Kim" 문자열 객체
\end{dhksmry}

이 상태에서 다음 코드가 실행된다.

\begin{dhkcode}
self.name = name
\end{dhkcode}

오른쪽의 \zsh{name}을 평가해 \zsh{"Kim"} 객체를 가져오고, \zsh{self}가 참조하는 인스턴스에 \zsh{name} 속성을 설정하는 것이다.

실행 후에는 다음과 같은 구조가 형성된다.

\begin{dhkdiagram}[in=6]
node p (1,1) :: p
node per (3,3) :: Person 인스턴스
node d (5,3) :: __dict__
node s (7,3) :: "name" -> "Kim"
edge p -> per :: HV
edge per -> d
edge d -> s
\end{dhkdiagram}

실제로 확인해 보자.

\begin{dhkcode}
class Person:
    def __init__(self, name):
        self.name = name

p = Person("Kim")

print(p.__dict__)
\end{dhkcode}

실행 결과:

\begin{dhkzsh}
{'name': 'Kim'}
\end{dhkzsh}

\begin{dhknote}
\begin{dhksmry}[centered]
매개변수 name과 인스턴스 속성 self.name은 다른 바인딩이다
\end{dhksmry}
\end{dhknote}

이 점이 특히 중요하다.

다음 코드를 보자.

\begin{dhkcode}
class Person:
    def __init__(self, name):
        self.name = name
        name = "Lee"

p = Person("Kim")

print(p.name)
\end{dhkcode}

실행 결과:

\begin{dhkzsh}
Kim
\end{dhkzsh}

왜 \zsh{"Lee"}가 출력되지 않을까?

\zsh{name="Lee"}는 지역 변수 \zsh{name}을 다른 문자열 객체에 재바인딩하는 코드이기 때문이다.

이미 설정한 \zsh{self.name} 속성을 변경하는 코드가 아니다.

반대로 다음처럼 작성하면 결과가 달라진다.

\begin{dhkcode}
class Person:
    def __init__(self, name):
        self.name = name
        self.name = "Lee"

p = Person("Kim")

print(p.name)
\end{dhkcode}

실행 결과:

\begin{dhkzsh}
Lee
\end{dhkzsh}

이번에는 \zsh{self.name} 속성 자체를 다시 설정했기 때문이다.

즉,

\begin{dhksmry}[centered]
name $\neq$ self.name이라는 바인딩
\end{dhksmry}

두 표현식이 같은 객체를 참조할 수는 있지만, 동일한 이름 저장 위치를 의미하지는 않는다.

\subsubsection{인스턴스마다 서로 다른 속성을 가질 수 있는 이유}

2.3에서 하나의 클래스로 여러 인스턴스를 생성할 수 있다는 사실을 배웠다.

\begin{dhkcode}
class Person:
    pass

p1 = Person()
p2 = Person()
\end{dhkcode}

이때 \zsh{p1}과 \zsh{p2}는 서로 다른 인스턴스 객체다.

따라서 일반적인 클래스에서는 각 인스턴스가 자신의 속성 저장 공간을 가질 수 있다.

\begin{dhkcode}
p1.name = "Kim"
p2.name = "Lee"

print(p1.__dict__)
print(p2.__dict__)
\end{dhkcode}

실행 결과:

\begin{dhkzsh}
{'name': 'Kim'}
{'name': 'Lee'}
\end{dhkzsh}

관계를 나타내면 다음과 같다.

\begin{dhkgraph}
node i1 (2,1) :: __dict__ // name → "Kim"
node i2 (2,3) :: __dict__ // name → "Lee"

node c (1,2) :: 클래스 객체 // Person

group p1(i1) [accent] :: p1 인스턴스
group p2(i2) [accent] :: p2 인스턴스

split c -> p1, p2

caption :: 두 인스턴스는 동일한 클래스에 속하지만 별도의 속성 저장 공간을 가진다.
\end{dhkgraph}

여기서 클래스 객체는 하나지만, 인스턴스는 두 개다.

그리고 각 인스턴스는 서로 다른 \zsh{__dict__}를 가지고 있다.

따라서 다음 코드도 가능하다.

\begin{dhkcode}
p1.age = 30
p2.city = "Seoul"

print(p1.__dict__)
print(p2.__dict__)
\end{dhkcode}

실행 결과:

\begin{dhkzsh}
{'name': 'Kim', 'age': 30}
{'name': 'Lee', 'city': 'Seoul'}
\end{dhkzsh}

이 예제에서 \zsh{p1}은 \zsh{age} 속성을 가지고 있지만 \zsh{city} 속성은 없다.

반면 \zsh{p2}는 \zsh{city} 속성을 가지고 있지만 \zsh{age} 속성은 없다.

즉, 동일한 클래스의 인스턴스라고 해서 반드시 동일한 인스턴스 속성을 가져야 하는 것은 아니다.

다만 실무에서는 객체의 상태를 예측하기 쉽도록 \zsh{__init__()}에서 필요한 인스턴스 속성을 일관되게 초기화하는 경우가 많다.

\subsubsection{인스턴스 속성은 언제 생성되는가?}

앞에서는 \zsh{__init__()} 내부에서 속성을 생성하였다.

\begin{dhkcode}
class Person:

    def __init__(self, name):
        self.name = name
\end{dhkcode}

하지만 Python에서는 반드시 \zsh{__init__()} 안에서만 인스턴스 속성을 생성해야 하는 것은 아니다.

다음 코드도 가능하다.

\begin{dhkcode}
class Person:
    pass

p = Person()

p.name = "Kim"
p.age = 30
\end{dhkcode}

\zsh{Person} 클래스에는 \zsh{name}이나 \zsh{age}를 정의하지 않았지만, 인스턴스에 속성을 설정할 수 있다.

왜냐하면 일반적인 Python 인스턴스는 동적으로 속성을 추가할 수 있기 때문이다.

다음 코드를 확인해 보자.

\begin{dhkcode}
print(p.__dict__)
\end{dhkcode}

실행 결과:

\begin{dhkzsh}
{'name': 'Kim', 'age': 30}
\end{dhkzsh}

따라서 속성의 생성 시점은 반드시 인스턴스 생성 시점과 일치하지 않는다.

다만 다음 두 상황은 구분해야 한다.

\begin{dhknote}
\begin{dhksmry}[centered]
인스턴스 생성
\end{dhksmry}
\end{dhknote}

\begin{dhkcode}
p = Person()
\end{dhkcode}

새로운 인스턴스 객체가 생성되는 과정이다.

\begin{dhknote}
\begin{dhksmry}[centered]
인스턴스 속성 설정
\end{dhksmry}
\end{dhknote}

\begin{dhkcode}
p.name = "Kim"
\end{dhkcode}

이미 존재하는 인스턴스에 속성을 설정하는 과정이다.

이 둘은 서로 다른 동작이다.

\begin{dhknote}
\begin{dhksmry}[centered]
속성이 없는 상태에서 읽으면?
\end{dhksmry}
\end{dhknote}

\begin{dhkcode}
class Person:
    pass

p = Person()

print(p.name)
\end{dhkcode}

이 코드는 오류를 발생시킨다.

\begin{dhkzsh}
AttributeError: 'Person' object has no attribute 'name'
\end{dhkzsh}

반대로 속성을 설정한 후에는 읽을 수 있다.

\begin{dhkcode}
p.name = "Kim"

print(p.name)
\end{dhkcode}

실행 결과:

\begin{dhkzsh}
Kim
\end{dhkzsh}

이처럼 속성은 객체가 생성되었다는 이유만으로 자동으로 존재하는 것이 아니다.

해당 속성이 실제로 설정되어 있거나, 클래스 등에서 속성 탐색을 통해 발견될 수 있어야 한다.

\subsubsection{Python은 \zsh{p.name}을 어떻게 처리하는가?}

이제 중요한 질문으로 넘어간다.

\begin{dhkcode}
p.name
\end{dhkcode}

Python은 이 표현식을 어떻게 처리할까?

단순히 \zsh{p.__dict__["name"]}을 실행한다고 생각할 수 있지만, 실제 속성 접근 규칙은 그보다 복잡하다.

일반적인 사용자 정의 인스턴스에서는 \zsh{object.__getattribute__()}가 기본적인 속성 접근 동작을 제공한다.

여기에는 인스턴스의 속성뿐 아니라 클래스 속성, 상속 관계, 디스크립터 등이 관여할 수 있다.

우선 가장 단순한 경우부터 살펴본다.

\begin{dhkcode}
class Person:
    pass

p = Person()
p.name = "Kim"

print(p.name)
\end{dhkcode}

이 예제에서는 인스턴스의 \zsh{__dict__}에 \zsh{"name"}이 존재하므로 해당 값을 읽을 수 있다.

개념적으로 다음과 같이 생각할 수 있다.

\begin{dhksmry}[centered]
p.name
\da
p 객체 확인
\da
속성 탐색
\da
"name" 속성 발견
\da
"Kim" 반환
\end{dhksmry}

하지만 클래스에도 같은 이름의 속성이 있다면 어떻게 될까?

\begin{dhkcode}
class Person:
    name = "Unknown"

p = Person()

print(p.name)
\end{dhkcode}

실행 결과:

\begin{dhkzsh}
Unknown
\end{dhkzsh}

이번에는 인스턴스에 직접 \zsh{name} 속성을 설정하지 않았다.

그런데도 \zsh{p.name}으로 값을 읽을 수 있다.

왜냐하면 Python이 인스턴스에서 속성을 찾지 못했을 때 클래스에서도 속성을 탐색할 수 있기 때문이다.

다음 코드를 확인해 보자.

\begin{dhkcode}
print(p.__dict__)
print(Person.__dict__["name"])
\end{dhkcode}

실행 결과:

\begin{dhkzsh}
{}
Unknown
\end{dhkzsh}

\subsubsection{클래스 속성과 인스턴스 속성이 같은 이름을 가지면?}

다음 코드를 보자.

\begin{dhkcode}
class Person:
    name = "Unknown"

p = Person()

print(p.name)
\end{dhkcode}

실행 결과:

\begin{dhkzsh}
Unknown
\end{dhkzsh}

여기까지는 클래스 속성에서 값을 읽었다.

그런데 다음 코드를 실행하면 어떻게 될까?

\begin{dhkcode}
p.name = "Kim"

print(p.name)
print(Person.name)
\end{dhkcode}

실행 결과:

\begin{dhkzsh}
Kim
Unknown
\end{dhkzsh}

왜 이런 결과가 발생할까?

\zsh{p.name = "Kim"}은 일반적인 속성 설정 규칙에 따라 인스턴스의 \zsh{name} 속성을 생성하기 때문이다.

클래스 속성 \zsh{Person.name}을 수정한 것이 아니다.

실제로 확인해 보자.

\begin{dhkcode}
print(p.__dict__)
print(Person.__dict__["name"])
\end{dhkcode}

실행 결과:

\begin{dhkzsh}
{'name': 'Kim'}
Unknown
\end{dhkzsh}

따라서 다음과 같은 구조가 형성된다.

\begin{dhkdir}[in=6]
Person 클래스 객체
> name -> "Unknown"
\end{dhkdir}

\begin{dhkdir}[in=6]
p 인스턴스 객체
> name -> "Kim"
\end{dhkdir}

이때 \zsh{p.name}은 인스턴스 속성 \zsh{"Kim"}을 반환하고, \zsh{Person.name}은 클래스 속성 \zsh{"Unknown"}을 반환한다.

이러한 현상을 일반적으로 

\zsh{인스턴스 속성이 클래스 속성을 가린다 (Shadowing)}고 설명한다.

다만 이 규칙에는 중요한 예외가 있다.

클래스에 같은 이름의 \zsh{데이터 디스크립터(Data Descriptor)}가 존재하면 인스턴스 \zsh{__dict__}보다 우선할 수 있다.

따라서 일반적인 속성 탐색 순서를 조금 더 정확하게 표현하면 다음과 같다.

\begin{dhkgraph}[uniform]
node i1 (1,1) [accent] :: p.name 속성 접근
node i2 (2,1):: 클래스 계층의 데이터 디스크립터
node i3 (3,1):: 인스턴스 __dict__
node i4 (4,1):: 클래스 계층의 일반 속성 / 비데이터 디스크립터
node i5 (5,1):: __getattr__ 등 후속 처리

edge i1 -> i2
edge i2 -> i3
edge i3 -> i4
edge i4 -> i5

caption :: 일반적인 인스턴스 속성 읽기의 개념적 우선순위. // 각 단계에서 속성이 발견되면 그에 맞는 처리가 이루어진다.
\end{dhkgraph}

이 도식은 기본적인 \zsh{object.__getattribute__()} 동작을 단순화한 것이다.

클래스에서 속성 접근을 직접 재정의하면 동작이 달라질 수 있다.

지금 단계에서는 인스턴스 속성과 클래스 속성이 별도의 저장 공간에 존재하며, Python은 정해진 규칙에 따라 속성을 탐색한다는 점을 이해하면 충분하다.

\subsubsection{\zsh{getattr(), setattr(), delattr()} - 속성을 동적으로 다루기}

지금까지는 점(.)을 사용해 속성에 접근했다.

\begin{dhkcode}
p.name
p.name = "Kim"
\end{dhkcode}

그런데 Python에서는 속성 이름을 문자열로 전달해 접근할 수도 있다.

이를 위해 다음 내장 함수를 제공한다.

\begin{dhktable}
|함수|	역할|
|:---|:---|
|getattr(obj, name)|	객체의 속성 읽기|
|setattr(obj, name, value)|	객체의 속성 설정|
|delattr(obj, name)|	객체의 속성 삭제|
|hasattr(obj, name)|	속성 접근 가능 여부 확인|
\end{dhktable}

예제를 살펴본다.

\begin{dhkcode}
class Person:
    pass

p = Person()

setattr(p, "name", "Kim")

print(getattr(p, "name"))
\end{dhkcode}

실행 결과:

\begin{dhkzsh}
Kim
\end{dhkzsh}

다음 두 코드는 일반적인 경우 같은 효과를 가진다.

\begin{dhkcode}
p.name = "kim"

setattr(p, "name", "Kim")
\end{dhkcode}

마찬가지로 다음 두 코드는 같은 속성을 읽는다.

\begin{dhkcode}
p.name

getattr(p, "name")
\end{dhkcode}

기본값을 지정할 수도 있다.

\begin{dhkcode}
print(getattr(p, "age", 0))
\end{dhkcode}

실행 결과:

\begin{dhkzsh}
0
\end{dhkzsh}

\zsh{age} 속성이 존재하지 않으므로 기본 값 0을 반환한다.

반면 다음 코드는 오류가 발생할 수 있다.

\begin{dhkcode}
print(getattr(p, "age"))
\end{dhkcode}

\zsh{getattr()}에 기본값을 전달하지 않았기 때문이다.

다음 코드는 일반적인 인스턴스에서 \zsh{name} 속성을 삭제한다.

\begin{dhkcode}
delattr(p, "name")

print(p.__dict__)
\end{dhkcode}

실행 결과:

\begin{dhkzsh}
{}
\end{dhkzsh}

단, \zsh{delattr()} 역시 Python의 속성 삭제 규칙을 사용하므로 항상 단순한 딕셔너리 삭제와 동일한 것은 아니다.

왜 이런 함수들이 필요할까?

속성 이름이 실행 중에 결정되는 상황에서 유용하기 때문이다.

\begin{dhkcode}
attribute_name = "name"

value = getattr(p, attribute_name, None)
\end{dhkcode}

여기서는 어떤 속성을 읽을지 문자열 변수로 지정할 수 있다.

이런 방식은 설정 객체, 데이터 변환, 프레임워크 내부 코드 등에서 활용된다.

\subsubsection{Mutable 객체를 인스턴스 속성에 저장하면?}

이번에는 이전에 학습한 Mutable / Immutable 개념과 연결해 보자.

\begin{dhkcode}
class Person:
    def __init__(self, name):
        self.name = name
        self.hobbies = []
\end{dhkcode}

여기서 \zsh{self.hobbies}에는 리스트 객체가 저장된다.

정확히는 인스턴스 속성이 리스트 객체를 참조하게 되는 것이다.

\begin{dhkcode}
p = Person("Kim")
p.hobbies.append("Reading")
\end{dhkcode}

이때 \zsh{append()}는 리스트 객체를 변경하는 메서드다.

따라서 다음 결과를 얻을 수 있다.

\begin{dhkcode}
print(p.hobbies)
\end{dhkcode}

실행 결과:

\begin{dhkzsh}
['Reading']
\end{dhkzsh}

그런데 다음 두 코드는 서로 다른 동작이다.

\begin{dhkcode}
p.hobbies.append("Music")

p.hobbies = ["Music"]
\end{dhkcode}

첫 번재 코드는 기존 리스트 객체를 변경한다.

두 번째 코드는 새로운 리스트 객체를 생성하고 \zsh{p.hobbies} 속성이 그 객체를 참조하도록 재설정한다.

이 차이를 \zsh{id()}로 확인할 수 있다.

\begin{dhkcode}
p = Person("Kim")

before = id(p.hobbies)

p.hobbies.append("Reading")

after_append = id(p.hobbies)

p.hobbies = ["Music"]

after_rebind = id(p.hobbies)

print(before == after_append)
print(before == after_rebind)
\end{dhkcode}

실행 결과:

\begin{dhkzsh}
True
False
\end{dhkzsh}

첫 번째는 기존 리스트 객체를 변경했기 때문에 객체의 동일성이 유지된다.

두 번째는 새 리스트를 만들어 속성을 재바인딩했기 때문에 다른 객체를 참조한다.

이것은 우리가 앞서 함수 인수 전달, 리스트 등을 공부하며 살펴본 원리와 같다.

\begin{dhknote}
\begin{dhksmry}[centered]
서로 다른 인스턴스의 리스트는 공유될까?
\end{dhksmry}
\end{dhknote}

\begin{dhkcode}
p1 = Person("Kim")
p2 = Person("Lee")

print(p1.hobbies is p2.hobbies)
\end{dhkcode}

실행 결과:

\begin{dhkzsh}
False
\end{dhkzsh}

왜냐하면 \zsh{__init__()}이 호출될 때마다 다음 표현식이 새 리스트 객체를 생성하기 때문이다.

\begin{dhkcode}
self.hobbies = []
\end{dhkcode}

따라서 각 인스턴스는 별도의 리스트를 참조한다.

반대로 클래스 속성으로 리스트를 정의하면 공유 문제가 발생할 수 있다.

\begin{dhkcode}
class Person:
    hobbies = []


p1 = Person()
p2 = Person()

p1.hobbies.append("Reading")

print(p2.hobbies)
\end{dhkcode}

실행 결과:

\begin{dhkzsh}
['Reading']
\end{dhkzsh}

두 인스턴스가 같은 클래스 속성의 리스트 객체를 참조하기 때문이다.

이 문제는 \zsh{2.7. 클래스 속성과 인스턴스 속성}에서 더 깊게 분석할 것이다.

\subsubsection{PyTorch와의 연결 - \zsh{self.linear}는 어디에 저장되는가?}

이제 PyTorch 모델 코드를 살펴보자.

\begin{dhkcode}
import torch.nn as nn


class MyModel(nn.Module):

    def __init__(self):
        super().__init__()

        self.linear = nn.Linear(10, 5)

    def forward(self, x):
        return self.linear(x)
\end{dhkcode}

여기서 다음 표현식에 주목하자.

\begin{dhkcode}
self.linear = nn.Linear(10, 5)
\end{dhkcode}

일반적인 Python 클래스였다면 \zsh{self.linear}는 인스턴스 속성으로 저장되는 경우가 많다.

하지만 PyTorch의 \zsh{nn.Module}에서는 조금 더 특별한 처리가 이루어진다.

\zsh{nn.Module}은 \zsh{__setattr__()}을 재정의하여, \zsh{nn.Module}의 하위 모듈이 인스턴스 속성으로 설정될 때 이를 등록하도록 구성되어 있다.

따라서 \zsh{self.linear}는 일반적인 인스턴스 딕셔너리에 단순히 저장되는 것이 아니라, PyTorch가 관리하는 하위 모듈 저장 구조에 등록된다.

실제로 확인해 보자.

\begin{dhkcode}
model = MyModel()

print(model._modules)
\end{dhkcode}

개념적인 출력:

\begin{dhkzsh}
{'linear': Linear(in_features=10, out_features=5, bias=True)}
\end{dhkzsh}

이때 \zsh{_modules}는 PyTorch가 하위 모듈을 관리하기 위해 사용하는 내부 구조다.

다음 코드를 통해 등록된 하위 모듈에 접근할 수 있다.

\begin{dhkcode}
print(model.linear)
\end{dhkcode}

출력:

\begin{dhkzsh}
Linear(in_features=10, out_features=5, bias=True)
\end{dhkzsh}

여기서 중요한 질문이 생긴다.

\zsh{model.linear}가 일반적인 \zsh{model.__dict__}에 직접 저장되지 않았는데도 왜 접근할 수 있을까?

그 이유는 \zsh{nn.Module}이 속성 접근과 설정을 위한 특수 메서드를 활용하기 때문이다.

일반적인 Python 객체에서는 속성 접근이 \zsh{__getattribute__()}를 통해 처리되고, 속성을 찾지 못한 경우 \zsh{__getattr__()}이 호출될 수 있다.

PyTorch의 \zsh{nn.Module}은 이러한 메커니즘을 이용해 등록된 하위 모듈과 파라미터에 접근할 수 있도록 구성되어 있다.

따라서 다음과 같은 관계가 형성된다.

\begin{dhksmry}[centered]
model.lienar
\da
Python 속성 접근
\da
nn.Module의 속성 관리 메커니즘
\da
등록된 하위 모듈 탐색
\da
nn.Linear 인스턴스 반환
\end{dhksmry}

이 구조 덕분에 PyTorch는 모델 내부의 하위 모듈과 파라미터를 체계적으로 관리할 수 있다.

예를 들어 다음 메서드들이 등록된 객체들을 활용한다.

\begin{dhkcode}
model.parameters()
model.state_dict()
model.to("cpu")
\end{dhkcode}

여기서 핵심은 다음과 같다.

\begin{dhksmry}[centered]
PyTorch의 self.linear도

Python의 속성 접근 문법을 사용하지만,

실제 저장과 탐색 과정은 

nn.Module이 정의한 객체 관리 규칙에 의해 확장된다.
\end{dhksmry}

따라서 Python의 기본적인 속성 접근 규칙을 이해하면 PyTorch의 모델 등록 구조도 이해하기 쉬워진다.

\subsubsection{핵심 정리}

이번 학습에서는 인스턴스 속성과 네임스페이스 관계를 살펴보았다.

\begin{dhknote}
\begin{dhksmry}[centered]
속성은 객체에 연결된 이름이다.
\end{dhksmry}
\end{dhknote}

\begin{dhkcode}
p.name
\end{dhkcode}

여기서 \zsh{name}은 \zsh{p} 객체의 속성 이름이다.

\begin{dhknote}
\begin{dhksmry}[centered]
일반적인 사용자 정의 인스턴스는 __dict__에 속성을 저장한다.
\end{dhksmry}
\end{dhknote}

\begin{dhkcode}
p.__dict__
\end{dhkcode}

인스턴스에 직접 설정된 일반적인 속성들을 확인할 수 있다.

\begin{dhknote}
\begin{dhksmry}[centered]
일반 변수와 인스턴스 속성은 다른 바인딩이다.
\end{dhksmry}
\end{dhknote}

\begin{dhkcode}
name = "Kim"

self.name = name
\end{dhkcode}

오른쪽 \zsh{name}은 지역 이름이고, 왼쪽 \zsh{self.name}은 인스턴스 속성이다.

\begin{dhknote}
\begin{dhksmry}[centered]
각 인스턴스는 독립적인 속성 저장 공간을 가질 수 있다.
\end{dhksmry}
\end{dhknote}

\begin{dhkcode}
p1.name = "Kim"
p2.name = "Lee"
\end{dhkcode}

두 인스턴스의 속성은 서로 독립적으로 설정할 수 있다.

\begin{dhknote}
\begin{dhksmry}[centered]
인스턴스 속성은 __init__() 밖에서도 설정할 수 있다.
\end{dhksmry}
\end{dhknote}

\begin{dhkcode}
p.age = 30
\end{dhkcode}

일반적인 Python 인스턴스는 동적으로 속성을 추가할 수 있다.

\begin{dhknote}
\begin{dhksmry}[centered]
속성 접근은 단순한 딕셔너리 조회보다 복잡하다.
\end{dhksmry}
\end{dhknote}

\begin{dhkcode}
p.name
\end{dhkcode}

Python은 클래스 속성, 디스크립터, 상속 관계 등을 고려해 속성을 탐색한다.

\begin{dhknote}
\begin{dhksmry}[centered]
getattr(), setattr(), delattr()로 속성을 동적으로 다룰 수 있다.
\end{dhksmry}
\end{dhknote}

\begin{dhkcode}
getattr(p, "name")
setattr(p, "age", 30)
delattr(p, "age")
\end{dhkcode}

\begin{dhknote}
\begin{dhksmry}[centered]
Mutable 객체를 속성에 저장할 때는 변경과 재바인딩을 구분해야 한다.
\end{dhksmry}
\end{dhknote}

\begin{dhkcode}
p.hobbies.append("Reading")
p.hobbies = ["Music"]
\end{dhkcode}

첫 번째는 기존 객체의 변경이고, 두 번째는 속성의 재바인딩이다.

\begin{dhknote}
\begin{dhksmry}[centered]
PyTorch는 Python의 속성 관리 메커니즘을 확장한다.
\end{dhksmry}
\end{dhknote}

\begin{dhkcode}
self.linear = nn.Linear(10, 5)
\end{dhkcode}

\zsh{nn.Module}은 하위 모듈을 등록하고 관리하기 위해 속성 설정 및 접근 메커니즘을 활용한다.

\subsubsection{2.4에서 반드시 이해해야 할 결론}

이번 학습에서 가장 중요한 관계는 다음과 같다.

\begin{dhksmry}[centered]
Instance Object \ra Attribute Namespace \ra Name-Object Binding
\end{dhksmry}

일반적인 Python 인스턴스에서는 이 구조를 다음처럼 확인할 수 있다.

\begin{dhkcode}
p.__dict__
\end{dhkcode}

그리고 다음 표현식은

\begin{dhkcode}
sefl.name = name
\end{dhkcode}

지역 이름 \zsh{name}이 참조하는 객체를, \zsh{self}가 참조하는 인스턴스의 \zsh{name} 속성에 연결하는 동작이다.

이것이 2.4의 핵심이다.